% ============================================================
%  CHAPITRE 3 — CONCEPTION DU SYSTÈME
% ============================================================
\chapter{Conception du système}

\section{Introduction}

Ce chapitre présente la conception architecturale complète du système de maintenance
prédictive. Nous détaillons les besoins fonctionnels et non fonctionnels, la modélisation
UML à travers cinq diagrammes, l'architecture microservices, le schéma de la base de
données, et les flux de données entre les composants.

La conception a été guidée par trois principes fondamentaux : la \textbf{séparation des
responsabilités} entre les services (chaque service a un rôle unique et bien défini),
la \textbf{résilience} face aux pannes réseau (retry exponentiel, verrous distribués),
et la \textbf{scalabilité} pour supporter un nombre croissant de machines surveillées.

% ─────────────────────────────────────────────────────────────
\section{Besoins fonctionnels}
% ─────────────────────────────────────────────────────────────

Les besoins fonctionnels ont été identifiés à partir de l'analyse des cas d'utilisation
et des entretiens avec les utilisateurs cibles (techniciens informatiques et administrateurs
système). Ils sont organisés par module fonctionnel :

\begin{table}[H]
\centering
\caption{Besoins fonctionnels du système}
\begin{tabularx}{\textwidth}{|c|l|X|}
\hline
\rowcolor{lightblue}
\textbf{ID} & \textbf{Fonctionnalité} & \textbf{Description} \\
\hline
BF01 & Collecte automatique & L'agent collecte CPU, RAM, disque et SMART toutes les heures \\
\hline
BF02 & Authentification & Connexion sécurisée par JWT avec rôles Admin/Technicien \\
\hline
BF03 & Prédiction ML & Calcul de la probabilité de panne à 7, 14 et 30 jours \\
\hline
BF04 & Alertes automatiques & Création d'alertes et envoi d'emails pour risques HIGH/CRITICAL \\
\hline
BF05 & Tableau de bord & Visualisation des métriques, prédictions RF (Proba. Panne) et LSTM (Santé Disque) côte à côte, et alertes \\
\hline
BF06 & Gestion des alertes & Accusé de réception et résolution des alertes \\
\hline
BF07 & Chatbot IA & Assistant conversationnel en langage naturel \\
\hline
BF08 & Gestion utilisateurs & Création, modification et suppression de comptes (Admin) \\
\hline
BF09 & Export CSV & Export des données de métriques au format CSV \\
\hline
BF10 & Détection d'anomalies & Isolation Forest pour détecter les comportements anormaux \\
\hline
BF11 & Prédiction LSTM & Prédiction de défaillance disque via le modèle LSTM Backblaze \\
\hline
BF12 & Versioning modèles & Registre des modèles ML avec activation conditionnelle \\
\hline
\end{tabularx}
\end{table}

% ─────────────────────────────────────────────────────────────
\section{Besoins non fonctionnels}
% ─────────────────────────────────────────────────────────────

\begin{table}[H]
\centering
\caption{Besoins non fonctionnels du système}
\begin{tabularx}{\textwidth}{|l|X|}
\hline
\rowcolor{lightblue}
\textbf{Critère} & \textbf{Exigence} \\
\hline
Performance & Temps de réponse API $<$ 100ms, requêtes DB $<$ 50ms \\
\hline
Disponibilité & Collecte automatique 24h/24, prédictions quotidiennes à 2h00 \\
\hline
Sécurité & JWT (24h), bcrypt (10 rounds), protection injection SQL via ORM Sequelize \\
\hline
Scalabilité & Architecture microservices permettant le déploiement indépendant de chaque service \\
\hline
Fiabilité & Retry exponentiel pour l'agent, verrou distribué PostgreSQL pour le scheduler \\
\hline
Confidentialité & LLM local (Ollama), aucune donnée transmise à des services cloud \\
\hline
Maintenabilité & Code modulaire, versioning des modèles ML, logs rotatifs \\
\hline
Portabilité & Conteneurisation Docker, compatible Windows et Linux \\
\hline
\end{tabularx}
\end{table}

% ─────────────────────────────────────────────────────────────
\section{Diagramme de classes}
% ─────────────────────────────────────────────────────────────

\subsection{Description}

Le diagramme de classes représente la structure statique du système. L'entité centrale
est \textbf{Machine}, qui est associée à plusieurs types de données collectées
(\texttt{SystemMetrics}, \texttt{SmartData}), aux résultats d'analyse (\texttt{Prediction},
\texttt{Anomaly}), et aux notifications générées (\texttt{Alert}). Les modèles ML sont
versionnés dans \texttt{MLModel}. Les utilisateurs sont gérés via \texttt{User} avec un
système de rôles (Admin / Viewer).

\subsection{Diagramme}

\begin{figure}[H]
  \centering
  \includegraphics[width=0.95\textwidth]{class_diagram.png}
  \caption{Diagramme de classes du système}
\end{figure}

\subsection{Interprétation des relations}

Ce diagramme illustre les relations suivantes :
\begin{itemize}
  \item Une \textbf{Machine} génère plusieurs \texttt{SystemMetrics}
        (relation 1..* — collecte horaire via l'agent) ;
  \item Une \textbf{Machine} possède plusieurs \texttt{SmartData}
        (données SMART du disque dur, collectées simultanément) ;
  \item Une \textbf{Machine} est associée à plusieurs \texttt{Prediction}
        (historique des prédictions ML générées chaque nuit à 2h00) ;
  \item Une \textbf{Machine} peut générer plusieurs \texttt{Alert}
        (alertes HIGH/CRITICAL avec statut ACTIVE/ACKNOWLEDGED/RESOLVED) ;
  \item Un \textbf{MLModel} est référencé par les \texttt{Prediction} via
        \texttt{model\_version}, permettant de tracer quelle version du modèle
        a produit chaque prédiction.
\end{itemize}

Les données collectées par l'agent (\texttt{SystemMetrics} et \texttt{SmartData}) sont
utilisées par les modèles de machine learning pour détecter les anomalies et générer les
prédictions. Si un problème est détecté (probabilité de panne $\geq 50\%$), le système
crée automatiquement une \texttt{Alert} et envoie une notification par email.

% ─────────────────────────────────────────────────────────────
\section{Diagramme de cas d'utilisation}
% ─────────────────────────────────────────────────────────────

\subsection{Description}

Le diagramme de cas d'utilisation identifie quatre acteurs dans le système :
\begin{itemize}
  \item \textbf{Administrateur} : gère les utilisateurs, les machines et a accès à toutes
        les fonctionnalités du système ;
  \item \textbf{Technicien} : surveille le système via le tableau de bord, gère les alertes,
        consulte les prédictions et interagit avec le chatbot ;
  \item \textbf{Agent de collecte} : programme Python autonome qui collecte et envoie les
        données automatiquement toutes les heures ;
  \item \textbf{Système ML} : service automatique qui génère les prédictions et déclenche
        les alertes sans intervention humaine.
\end{itemize}

\subsection{Diagramme}

\begin{figure}[H]
  \centering
  \includegraphics[width=0.95\textwidth]{use_case_diagram.png}
  \caption{Diagramme de cas d'utilisation}
\end{figure}

\subsection{Interprétation}

L'administrateur dispose de tous les droits, y compris la gestion des utilisateurs
(création, modification de rôle, suppression via le middleware \texttt{requireAdmin}).
Le technicien a un accès fonctionnel : il peut consulter le tableau de bord, visualiser
les détails des machines, gérer les alertes (accusé de réception, résolution) et interagir
avec le chatbot. L'agent de collecte et le système ML sont des acteurs automatiques qui
n'interagissent pas directement avec l'interface utilisateur.

% ─────────────────────────────────────────────────────────────
\section{Diagramme de séquence — Affichage du tableau de bord}
% ─────────────────────────────────────────────────────────────

\subsection{Description}

Ce diagramme illustre le flux d'interactions lors de l'ouverture du tableau de bord par
un technicien. L'utilisateur s'authentifie, le frontend envoie une requête au backend,
qui interroge la base de données et retourne les données agrégées (KPIs, liste des
machines, alertes actives).

\subsection{Diagramme}

\begin{figure}[H]
  \centering
  \includegraphics[width=0.95\textwidth]{sequence_dashboard.png}
  \caption{Diagramme de séquence — Affichage du tableau de bord}
\end{figure}

\subsection{Interprétation}

Le flux commence par l'authentification de l'utilisateur qui reçoit un token JWT signé
avec \texttt{JWT\_SECRET}, valide 24 heures. Ce token est inclus dans toutes les requêtes
suivantes dans l'en-tête \texttt{Authorization: Bearer <token>}. Le backend effectue deux
requêtes SQL principales : une pour les KPIs globaux (nombre de machines, alertes actives,
risque moyen) via \texttt{GET /api/dashboard/overview}, et une pour la liste complète des
machines avec leurs dernières prédictions RF et LSTM via \texttt{GET /api/dashboard/machines}.
Le frontend React reçoit les données JSON et effectue le rendu des composants (\texttt{KPICards},
\texttt{MachineList}, \texttt{AlertsList}), le composant \texttt{MachineList} affichant
notamment deux colonnes de prédiction distinctes : ``Proba. Panne (RF)'' et ``Santé Disque (LSTM)''.
Le temps de réponse total est inférieur à 100ms.

% ─────────────────────────────────────────────────────────────
\section{Diagramme de séquence — Génération de prédictions}
% ─────────────────────────────────────────────────────────────

\subsection{Description}

Ce diagramme illustre le pipeline de prédiction automatique déclenché chaque nuit à 2h00.
Le scheduler APScheduler réveille le service ML, qui extrait les caractéristiques, applique
le modèle Random Forest, calcule le niveau de risque, et génère les alertes si nécessaire.

\subsection{Diagramme}

\begin{figure}[H]
  \centering
  \includegraphics[width=0.95\textwidth]{sequence_prediction_alert.png}
  \caption{Diagramme de séquence — Génération de prédictions et alertes}
\end{figure}

\subsection{Interprétation}

Le pipeline commence par l'acquisition d'un verrou distribué PostgreSQL
(\texttt{pg\_try\_advisory\_lock(12345)}) pour éviter les exécutions concurrentes en cas
de redémarrage du service. Le service ML extrait ensuite 63 caractéristiques statistiques
sur les 30 derniers jours pour chaque machine. Le modèle Random Forest actif (récupéré
depuis le registre \texttt{ml\_models}) calcule la probabilité de panne. Si le risque est
HIGH ($\geq 50\%$) ou CRITICAL ($\geq 70\%$), une alerte est créée dans la base de données
et un email HTML est envoyé via SMTP (Nodemailer). Le verrou est libéré à la fin du
traitement. En cas d'échec, un retry est planifié 1 heure plus tard.

% ─────────────────────────────────────────────────────────────
\section{Diagramme de séquence — Collecte de données}
% ─────────────────────────────────────────────────────────────

\subsection{Description}

Ce diagramme illustre le cycle de collecte horaire de l'agent Python. L'agent collecte
les métriques système et SMART, construit le payload JSON, et l'envoie au backend avec
un mécanisme de retry en cas d'échec.

\subsection{Diagramme}

\begin{figure}[H]
  \centering
  \includegraphics[width=0.95\textwidth]{sequence_data_collection.png}
  \caption{Diagramme de séquence — Collecte de données par l'agent}
\end{figure}

\subsection{Interprétation}

L'agent exécute un cycle de collecte complet toutes les heures. Il collecte
séquentiellement les métriques CPU (usage + température), RAM (usage + disponible),
disque (usage + libre), puis les données SMART via \texttt{smartctl}. Si SMART n'est
pas disponible sur le matériel, des valeurs par défaut sont utilisées. Le payload JSON
est envoyé au backend avec un token Bearer. En cas d'erreur serveur (5xx), l'agent
réessaie avec un backoff exponentiel (1s, 2s, 4s). Les erreurs client (4xx) ne
déclenchent pas de retry.

% ─────────────────────────────────────────────────────────────
\section{Architecture générale du système}
% ─────────────────────────────────────────────────────────────

\subsection{Vue d'ensemble microservices}

Le système adopte une architecture microservices composée de quatre services indépendants.
Chaque service peut être démarré, arrêté et mis à l'échelle indépendamment des autres.
La communication inter-services se fait exclusivement via des API REST authentifiées.

\begin{figure}[H]
\centering
\includegraphics[width=0.95\textwidth]{architecture_microservices.png}
\caption{Architecture microservices du système PC Technician Assistant}
\end{figure}

\subsection{Responsabilités de chaque service}

\begin{table}[H]
\centering
\caption{Responsabilités des services de l'architecture microservices}
\begin{tabularx}{\textwidth}{|l|l|X|}
\hline
\rowcolor{lightblue}
\textbf{Service} & \textbf{Technologie} & \textbf{Responsabilité} \\
\hline
Agent & Python 3.9 & Collecte horaire des métriques système et SMART, envoi au backend \\
\hline
Backend & Node.js 20 + Express & API REST centrale, authentification JWT, routage, emails \\
\hline
ML Service & Python + Flask & Extraction de features, entraînement RF, inférence LSTM, scheduler \\
\hline
Frontend & React 18 + Vite & Interface utilisateur SPA, visualisation, chatbot \\
\hline
PostgreSQL & v14 & Stockage persistant de toutes les données (9 tables) \\
\hline
Ollama & LLaMA 3.2:1b & Inférence LLM locale pour le chatbot \\
\hline
\end{tabularx}
\end{table}

\subsection{Flux de données complet}

\begin{enumerate}
  \item \textbf{Collecte (Agent $\rightarrow$ Backend $\rightarrow$ DB)} : toutes les heures,
        l'agent envoie les métriques au backend qui les persiste dans PostgreSQL via
        Sequelize ORM ;
  \item \textbf{Prédiction (DB $\rightarrow$ ML $\rightarrow$ DB)} : à 2h00, le service ML
        lit les données des 30 derniers jours, extrait 63 features, génère les prédictions
        et les stocke dans la table \texttt{predictions} ;
  \item \textbf{Alertes (ML $\rightarrow$ Backend $\rightarrow$ Email)} : pour les risques
        élevés, le backend crée une alerte dans \texttt{alerts} et envoie un email HTML
        via Nodemailer/SMTP ;
  \item \textbf{Visualisation (DB $\rightarrow$ Backend $\rightarrow$ Frontend)} : le
        technicien consulte le dashboard qui affiche les données via des requêtes REST
        authentifiées par JWT ;
  \item \textbf{Chatbot (Frontend $\rightarrow$ Backend $\rightarrow$ DB/Ollama)} : les
        questions sont traitées par détection d'intention, requêtes SQL déterministes
        ou LLM local selon le type de question.
\end{enumerate}

% ─────────────────────────────────────────────────────────────
\section{Schéma de la base de données}
% ─────────────────────────────────────────────────────────────

\subsection{Organisation des tables}

La base de données PostgreSQL est composée de 9 tables organisées autour de l'entité
centrale \texttt{machines}. Les migrations SQL sont versionnées de 001 à 010 et
exécutées dans l'ordre par le script \texttt{migrate.js}.

\begin{table}[H]
\centering
\caption{Description des tables de la base de données}
\begin{tabularx}{\textwidth}{|l|c|X|}
\hline
\rowcolor{lightblue}
\textbf{Table} & \textbf{Migration} & \textbf{Description} \\
\hline
\texttt{machines}        & 001 & Informations d'identification (hostname, IP, OS, serial) \\
\hline
\texttt{agents}          & 002 & Tokens Bearer d'authentification des agents de collecte \\
\hline
\texttt{system\_metrics} & 003 & Métriques horaires (cpu\_usage, memory\_usage, disk\_usage) \\
\hline
\texttt{smart\_data}     & 004 & Données SMART (health\_status, temperature, read/write\_errors) \\
\hline
\texttt{predictions}     & 005 & Probabilités de panne à 7/14/30 jours + contributing\_factors JSON \\
\hline
\texttt{anomalies}       & 006 & Anomalies détectées par Isolation Forest \\
\hline
\texttt{ml\_models}      & 007--010 & Registre des modèles RF et LSTM avec versioning, métriques et statut actif \\
\hline
\texttt{alerts}          & 008 & Alertes avec statut ACTIVE/ACKNOWLEDGED/RESOLVED \\
\hline
\texttt{users}           & 009 & Comptes utilisateurs avec rôles et mots de passe bcrypt \\
\hline
\end{tabularx}
\end{table}

\subsection{Colonnes clés de la table \texttt{predictions}}

La table \texttt{predictions} stocke les résultats du pipeline ML avec les colonnes
suivantes :
\begin{itemize}
  \item \texttt{failure\_probability\_7d}, \texttt{failure\_probability\_14d},
        \texttt{failure\_probability\_30d} : probabilités en pourcentage (0--100) ;
  \item \texttt{risk\_level} : énumération LOW/MEDIUM/HIGH/CRITICAL ;
  \item \texttt{model\_version} : référence à l'ID du modèle dans \texttt{ml\_models},
        qui supporte les types \texttt{random\_forest}, \texttt{isolation\_forest} et
        \texttt{lstm} (ajouté par la migration 010) ;
  \item \texttt{contributing\_factors} : colonne JSONB stockant les 5 features les plus
        importantes avec leur valeur et leur score d'importance, affichées dans le
        composant \texttt{MachineDetails} du frontend.
\end{itemize}

\subsection{Verrou distribué PostgreSQL}

Pour éviter les exécutions concurrentes du scheduler de prédictions (par exemple lors
d'un redémarrage du service ML), le pipeline utilise le mécanisme de verrous consultatifs
de PostgreSQL :

\begin{itemize}
  \item \texttt{pg\_try\_advisory\_lock(12345)} : tente d'acquérir le verrou de manière
        non bloquante. Retourne \texttt{true} si le verrou est acquis, \texttt{false}
        si un autre processus le détient déjà ;
  \item \texttt{pg\_advisory\_unlock(12345)} : libère le verrou à la fin du traitement,
        même en cas d'exception (bloc \texttt{finally}).
\end{itemize}

% ─────────────────────────────────────────────────────────────
\section{Conception du système d'authentification}
% ─────────────────────────────────────────────────────────────

\subsection{Authentification multi-niveaux}

Le backend implémente trois mécanismes d'authentification distincts adaptés à chaque
type d'acteur :

\begin{enumerate}
  \item \textbf{Token statique Bearer pour les agents} : le token est vérifié dans la
        table \texttt{agents} de la base de données. Simple et rapide, sans gestion de
        session. Chaque machine dispose de son propre token unique.

  \item \textbf{JWT pour les utilisateurs} : token signé avec \texttt{JWT\_SECRET},
        expiration 24h, payload \texttt{\{id, email, role\}}. Stateless — le serveur
        n'a pas besoin de stocker les sessions.

  \item \textbf{Middleware RBAC \texttt{requireAdmin}} : vérifie que
        \texttt{req.user.role === 'admin'} pour les routes de gestion des utilisateurs.
        Doit être utilisé après \texttt{authenticateToken}.
\end{enumerate}

Les mots de passe sont hashés avec bcrypt (10 rounds de salt) avant stockage. La
comparaison utilise \texttt{bcrypt.compare()} qui est résistante aux attaques par timing.

% ─────────────────────────────────────────────────────────────
\section{Conclusion}
% ─────────────────────────────────────────────────────────────

Ce chapitre a présenté la conception complète du système à travers cinq diagrammes UML,
les besoins fonctionnels et non fonctionnels, l'architecture microservices, et le schéma
de la base de données. La modélisation UML confirme la cohérence de l'architecture et la
clarté des flux de données entre les composants. Le chapitre suivant détaille
l'implémentation technique de chaque module.
